%% toc: P15 172 172.3,172.4
% Proofs for the general-score bias bound, the asymmetric variance correction, the standing of the unknown-scale proposals and the pair-average median risk.
\section{Huber's general bias bound, proposal standing and pair-median risk: proofs}
\label{sec:annex-p15-w02b-p15-071-huber-completions}

Use the notation of
Chapter~\ref{sec:p15-w02b-p15-071-huber-completions}. Two elementary
facts recur. First, for a nondecreasing \(\psi\), each map
\(t\mapsto\psi(x-t)\) is nonincreasing, so every \(\lambda_K\) is
nonincreasing and any derivative of it is nonpositive. Second, if
\(\sqrt nT_n\) converges in law to \(L\), then for each \(C>0\) the
function \(x\mapsto\min(x^2,C^2)\) is bounded and continuous, so
\(E\min(nT_n^2,C^2)\to E\min(L^2,C^2)\); letting \(C\to\infty\),
monotone convergence gives \(EL^2\). Thus the limiting truncated risk is
the second moment of the limit law.

\begin{annexitem}{annex:p15-w02b-p15-071-bias-identity}{thm:p15-w02b-p15-071-bias-identity}{The mean-value bias identity and bound (2)}
\begin{proof}
(i) Skew symmetry at zero gives \(\psi(0)=-\psi(0)\), so
\(\psi(0)=0\). Since \(X\) and \(-X\) have the same law under \(G\),
\(\lambda_G(0)=E_G\psi(X)=E_G\psi(-X)=-E_G\psi(X)\), hence
\(\lambda_G(0)=0\). More generally
\(\lambda_G(-t)=E_G\psi(X+t)=E_G\psi(-X+t)=-\lambda_G(t)\), so
\(\lambda_G\) is odd. Linearity of expectation gives
\(\lambda_F=(1-\varepsilon)\lambda_G+\varepsilon\lambda_H\), hence
\(\lambda_F(0)=\varepsilon\lambda_H(0)\), and
\(|\lambda_H(0)|\leq E_H|\psi(X)|\leq s\).

(ii) Then \(\lambda_F\) is differentiable on the stated open interval,
in particular continuous on the closed segment. If \(c\neq0\), the
mean-value theorem gives \(\vartheta\in(0,1)\) with
\(\lambda_F(c)-\lambda_F(0)=c\lambda_F'(\vartheta c)\); as
\(\lambda_F(c)=0\), this is \(\lambda_F(0)+c\lambda_F'(\vartheta c)=0\).
If \(c=0\), then \(\lambda_F(0)=0\) and every \(\vartheta\) serves.
Write \(D=(1-\varepsilon)\lambda_G'(\vartheta c)
+\varepsilon\lambda_H'(\vartheta c)=\lambda_F'(\vartheta c)\).
By the first preliminary fact \(\lambda_H'(\vartheta c)\leq0\), so
\(\lambda_G'(\vartheta c)<0\) gives
\(D\leq(1-\varepsilon)\lambda_G'(\vartheta c)<0\).
Dividing, \(c=-\lambda_F(0)/D=-\varepsilon\lambda_H(0)/D\), and
\[
 |c|=\frac{\varepsilon|\lambda_H(0)|}{|D|}\leq
 \frac{\varepsilon s}{(1-\varepsilon)(-\lambda_G'(\vartheta c))}.
\]

(iii) From \(\lambda_F(c)=0\),
\((1-\varepsilon)\lambda_G(c)=-\varepsilon\lambda_H(c)\).
If \(c\neq0\), the mean-value theorem and \(\lambda_G(0)=0\) give
\(\vartheta'\in(0,1)\) with \(\lambda_G(c)=c\lambda_G'(\vartheta'c)\);
if \(c=0\) any \(\vartheta'\) serves. Hence
\((1-\varepsilon)c\lambda_G'(\vartheta'c)=-\varepsilon\lambda_H(c)\).
Since \(|\lambda_H(c)|\leq s\) and \(\lambda_G'(\vartheta'c)<0\),
dividing by \((1-\varepsilon)|\lambda_G'(\vartheta'c)|\) gives the bound.
No property of \(\lambda_H\) other than boundedness was used.

(iv) Let \(\lambda_G\) be differentiable on \((-\tau_0,\tau_0)\).
For \(0<\tau<\tau_0\) put \(\mu(\tau)=-\lambda_G(\tau)=\lambda_G(-\tau)\),
using oddness. Since \(\lambda_G(u)/u\to\lambda_G'(0)<0\) as
\(u\downarrow0\), \(\lambda_G(u)<0\) for all sufficiently small \(u>0\);
by monotonicity \(\lambda_G(\tau)\leq\lambda_G(u)<0\) for such \(u\leq\tau\).
Hence \(\mu(\tau)>0\).
Let \(c\) be a zero of \(\lambda_F\) for some \(H\). If \(c\geq\tau\),
then \(\lambda_G(c)\leq\lambda_G(\tau)=-\mu(\tau)\); if \(c\leq-\tau\),
then \(\lambda_G(c)\geq\mu(\tau)\). In both cases
\(|\lambda_G(c)|\geq\mu(\tau)\). But
\(|\lambda_G(c)|=\varepsilon|\lambda_H(c)|/(1-\varepsilon)
\leq\varepsilon s/(1-\varepsilon)\). Consequently, whenever
\(\varepsilon s/(1-\varepsilon)<\mu(\tau)\), every zero for every \(H\)
satisfies \(|c|<\tau\), so \(\overline c_\varepsilon\leq\tau\).
As \(\varepsilon\downarrow0\) the left side tends to zero for fixed
\(\tau\); since \(\tau\) is arbitrary, \(\overline c_\varepsilon\to0\).

Put \(\omega(\tau)=\sup_{|u|\leq\tau}|\lambda_G'(u)-\lambda_G'(0)|\),
which tends to zero by continuity of \(\lambda_G'\) at zero.
For \(\varepsilon\) small enough that \(\overline c_\varepsilon<\tau_0\)
and \(\omega(\overline c_\varepsilon)<-\lambda_G'(0)\), part (iii)
applies to every zero, with \(|\vartheta'c|\leq\overline c_\varepsilon\),
so \(-\lambda_G'(\vartheta'c)\geq-\lambda_G'(0)-\omega(\overline c_\varepsilon)>0\).
Taking the supremum over \(H\) and zeros,
\[
 \overline c_\varepsilon\leq
 \frac{\varepsilon s}{(1-\varepsilon)(-\lambda_G'(0))}\cdot
 \left(1-\frac{\omega(\overline c_\varepsilon)}{-\lambda_G'(0)}\right)^{-1},
\]
and the last factor tends to one.
\end{proof}
\end{annexitem}

\begin{annexitem}{annex:p15-w02b-p15-071-bias-variance}{thm:p15-w02b-p15-071-bias-variance}{Admitting asymmetric contaminants changes the worst variance by o(c)}
\begin{proof}
\emph{The criterion.} Since \(\psi\) is continuous and nondecreasing,
\(\rho(u)=\int_0^u\psi\) is continuous, convex and has derivative
\(\psi\), so its left and right derivatives coincide with \(\psi\).
As \(\psi\) is odd, nondecreasing and not identically zero, some
\(u_0>0\) has \(\psi(u_0)>0\); then \(\rho(u)\geq\psi(u_0)(u-u_0)\) for
\(u\geq u_0\), and \(\rho\) is even, so \(\rho(u)\to\infty\) as
\(u\to\pm\infty\). Thus \cref{def:p15-convex-location-criterion}
applies with this \(\psi\).

(i) Let \(H\in\mathcal H_\varepsilon\) with zero \(c\). Then
\(\lambda_F'(c)=(1-\varepsilon)\lambda_G'(c)+\varepsilon\lambda_H'(c)
\leq-(1-\varepsilon)b_G(c)<0\), since \(\lambda_H'(c)\leq0\).
By the definition of the derivative there is \(\delta>0\) with
\(\lambda_F<0\) on \((c,c+\delta)\) and \(\lambda_F>0\) on
\((c-\delta,c)\); monotonicity extends these signs to all \(t>c\) and
all \(t<c\). So the zero is unique and
\cref{thm:p15-convex-location-consistency} gives consistency for \(c\).
The function \(t\mapsto E_F\psi(X-t)^2\) is bounded by \(s^2\) and is
continuous by dominated convergence, because \(\psi\) is continuous and
bounded. \Cref{thm:p15-convex-location-normality}, with
\(a=-\lambda_F'(c)>0\), gives the normal limit with variance
\(E_F\psi(X-c)^2/\lambda_F'(c)^2\). Its numerator is
\((1-\varepsilon)a_G(c)+\varepsilon E_H\psi(X-c)^2\leq
(1-\varepsilon)a_G(c)+\varepsilon s^2\), and its denominator is at least
\((1-\varepsilon)^2b_G(c)^2\).

(ii) Replace \(N\) by \(N\cap(-N)\), still a neighborhood of zero.
By skew symmetry and the symmetry of \(G\),
\(a_G(-t)=E_G\psi(X+t)^2=E_G\psi(-X+t)^2=E_G\psi(X-t)^2=a_G(t)\).
The proof of \cref{thm:p15-w02b-p15-071-bias-identity}(i) shows
\(\lambda_G\) is odd; differentiating \(\lambda_G(-t)=-\lambda_G(t)\)
gives \(\lambda_G'(-t)=\lambda_G'(t)\), so \(b_G\) is even.
An even function \(e\) differentiable at zero has
\(e'(0)=\lim_{h\to0}(e(h)-e(0))/h=\lim_{h\to0}(e(-h)-e(0))/h=-e'(0)\),
so \(e'(0)=0\).
Since \(b_G\) is differentiable, hence continuous, at zero and
\(b_G(0)>0\), there is \(\tau_1>0\) with \(b_G>0\) on
\([-\tau_1,\tau_1]\subseteq N\). There put \(u=a_G/b_G^2\), \(v=1/b_G^2\),
both even and differentiable at zero, so
\(u(t)-u(0)=o(t)\) and \(v(t)-v(0)=o(t)\). Then
\[
 R_\varepsilon(t)=\frac{u(t)}{1-\varepsilon}
           +\frac{\varepsilon s^2}{(1-\varepsilon)^2}v(t).
\]
For \(0\leq\varepsilon\leq1/2\) both coefficients are at most two.
For \(0<\tau\leq\tau_1\) set
\[
 \eta(\tau)=\sup_{0<|t|\leq\tau}
 \frac{2|u(t)-u(0)|+2s^2|v(t)-v(0)|}{|t|},
\]
which is finite for small \(\tau\) and tends to zero as \(\tau\downarrow0\).
Hence \(R_\varepsilon(t)-R_\varepsilon(0)\leq|t|\eta(\tau)\leq\tau\eta(\tau)\)
for \(|t|\leq\tau\).
If \(\lambda_G'\) is continuous at zero, then since
\(\lambda_G'(0)=-b_G(0)<0\), part (iv) of
\cref{thm:p15-w02b-p15-071-bias-identity} gives
\(\overline c_\varepsilon\to0\). For \(\varepsilon\leq1/2\) small enough
that \(\overline c_\varepsilon\leq\tau_1\), every \(H\in\mathcal H_\varepsilon\)
has its zero in \([-\overline c_\varepsilon,\overline c_\varepsilon]\),
and part (i) gives
\(V(\psi,F)\leq R_\varepsilon(c)\leq
R_\varepsilon(0)+\overline c_\varepsilon\eta(\overline c_\varepsilon)\).
As \(\eta(\overline c_\varepsilon)\to0\), the correction is
\(o(\overline c_\varepsilon)\).

(iii) For each \(x\), \(t\mapsto\psi_k(x-t)\) is Lipschitz with
constant one and has derivative \(-{\bf1}_{\{|x-t|<k\}}\) except where
\(|x-t|=k\). For a law \(K\) with \(K\{t-k,t+k\}=0\), the difference
quotients of \(\psi_k(X-\cdot)\) at \(t\) are bounded by one and converge
\(K\)-almost surely, so dominated convergence gives
\(\lambda_K'(t)=-K((t-k,t+k))\). For \(K=\Phi\) this holds at every \(t\),
so \(N=\mathbb R\) and \(b_G(t)=\Phi(t+k)-\Phi(t-k)>0\), with derivative
\(\phi(t+k)-\phi(t-k)\). Also
\(\partial_t\min((z-t)^2,k^2)=-2\psi_k(z-t)\) except at \(|z-t|=k\),
with bound \(2k\), so \(a_G'(t)=-2\lambda_\Phi(t)\) exists everywhere.
Both \(a_G'\) and \(b_G'\) are thus available at zero, and \(\lambda_\Phi'\)
is continuous. By \cref{thm:p15-huber-asymmetric-bias}, \(\lambda_F\) has a
unique zero \(c\) for every \(H\); if \(H\{c-k,c+k\}=0\), the formula
just proved shows \(\lambda_H\) is differentiable at \(c\), so
\(H\in\mathcal H_\varepsilon\). Further
\(a_G(0)=E\min(Z^2,k^2)=\int_{-k}^kz^2\phi(z)\,dz+2k^2(1-\Phi(k))
=A(k)-2k\phi(k)+2k^2(1-\Phi(k))=B(k)\), using
\(\int_{-k}^kz^2\phi=A(k)-2k\phi(k)\) by parts, and \(b_G(0)=A(k)\).
Hence \(R_\varepsilon(0)=V^*(k,\varepsilon)\).
A symmetric \(H\) supported on \(\{|x|>k\}\) has zero \(c=0\),
no mass at \(\pm k\), and variance exactly \(V^*(k,\varepsilon)\) by
\cref{thm:p15-huber-normal-variance}; this gives the lower inequality.
The upper inequality is part (ii).
\end{proof}
\end{annexitem}

\begin{annexitem}{annex:p15-w02b-p15-071-winsor}{thm:p15-w02b-p15-071-winsor}{Each proposal's location is a Winsorized mean}
\begin{proof}
Let \(\tilde x_i\) be the modified observations. If
\(|x_i-T|\leq kS\), then \(\psi(k,(x_i-T)/S)=(x_i-T)/S\) and
\(\tilde x_i=x_i\). If \(x_i>T+kS\), then \(\psi(k,(x_i-T)/S)=k\) and
\(\tilde x_i-T=kS\). If \(x_i<T-kS\), then the score is \(-k\) and
\(\tilde x_i-T=-kS\). In every case
\(S\psi(k,(x_i-T)/S)=\tilde x_i-T\). Summing and using the equation,
\(\sum_i(\tilde x_i-T)=0\), so \(T=n^{-1}\sum_i\tilde x_i\).
The location equations of Proposals 1 and 2 have this form, the
factor \(n^{-1}\) being immaterial, and Proposal 3's side condition is
the case \(S=1\). The description of the modified set is the case
analysis just given.
\end{proof}
\end{annexitem}

\begin{annexitem}{annex:p15-w02b-p15-071-prop2-minimax}{thm:p15-w02b-p15-071-prop2-minimax}{Proposal 2 with matched clipping attains the minimax value}
\begin{proof}
\emph{The matched constant.} For \(\ell>0\),
\(B(\ell)/\ell^2=E\min(Z^2/\ell^2,1)\) with \(Z\) standard normal,
since \(B(\ell)=E\min(Z^2,\ell^2)\) as computed in the proof of
\cref{thm:p15-w02b-p15-071-bias-variance}(iii). For \(\ell<\ell'\),
\(\min(z^2/\ell^2,1)\geq\min(z^2/\ell'^2,1)\), strictly on
\(\{0<|z|<\ell'\}\), which has positive probability; so the function is
strictly decreasing. It is continuous by dominated convergence, tends to
one as \(\ell\downarrow0\) because \(Z\neq0\) almost surely, and tends to
zero as \(\ell\to\infty\). Hence each value in \((0,1)\) is attained at
exactly one \(\ell\). Now
\(r_\varepsilon-\varepsilon=(1-\varepsilon)B(k_\varepsilon)/k_\varepsilon^2>0\)
and \(r_\varepsilon<\varepsilon+(1-\varepsilon)=1\), so a unique \(k\)
exists and \(\varepsilon<r_\varepsilon\).
The defining equation of \(q_{\max}\) in
\cref{thm:p15-huber-joint-variance-bound} is
\(B(q)/q^2=(r_\varepsilon-\varepsilon)/(1-\varepsilon)
=B(k_\varepsilon)/k_\varepsilon^2\), whose unique root is
\(q_{\max}=k_\varepsilon\). Then
\(W(k_\varepsilon,\varepsilon)=V^*(k_\varepsilon,\varepsilon)=1/I_*\)
by \cref{thm:p15-huber-least-favorable-normal}.

\emph{Risk.} The normal component has no atoms, so
\(F\{0\}=\varepsilon H\{0\}<1-r_\varepsilon=1-\beta/k^2\): this is the
strict identification condition of \cref{thm:p15-huber-joint-identified}.
By \cref{thm:p15-huber-joint-law-and-variance} the limiting truncated
risk exists and equals \(S_0^2EH_1^2\), and by
\cref{thm:p15-huber-joint-atomic-bound}, whose hypotheses
\(\beta=B(k)\), \(r=B(k)/k^2\), \(0\leq\varepsilon<r\) hold,
\(S_0^2EH_1^2\leq W(q_{\max},\varepsilon)=1/I_*\), with equality for
symmetric contaminants strictly outside \([-q_{\max},q_{\max}]\).
If \(\varepsilon<1-r_\varepsilon\), then every symmetric \(H\) satisfies
\(\varepsilon H\{0\}\leq\varepsilon<1-r_\varepsilon\). The limits exist,
so the \(\liminf\) in \(\mathcal R_\varepsilon\) is the limit and
\(\mathcal R_\varepsilon(T)=\sup_H S_0^2EH_1^2=1/I_*\). The equality with
the infimum over translation-equivariant estimators is
\cref{thm:p15-huber-location-minimax}.

\emph{Other constants.} If \(B(k)/k^2>\varepsilon\), the same theorem
defines \(q_{\max}\), and \(W(q_{\max},\varepsilon)=V^*(q_{\max},\varepsilon)\)
is, by \cref{thm:p15-huber-normal-variance}, the supremum over symmetric
\(H\) of the limiting variances of the known-scale clipped estimate with
cutoff \(q_{\max}\); by the preliminary fact this is its value of
\(\mathcal R_\varepsilon\). That estimate is translation equivariant by
\cref{thm:p15-convex-location-minimizers}, so
\cref{thm:p15-huber-location-minimax} gives
\(W(q_{\max},\varepsilon)\geq1/I_*\).
The numerical table is an evaluation of these equations.
\end{proof}
\end{annexitem}

\begin{annexitem}{annex:p15-w02b-p15-071-prop1-population}{thm:p15-w02b-p15-071-prop1-population}{Proposal 1: the population worst case}
\begin{proof}
The map \(q\mapsto F((-q,q))=(1-\varepsilon)A(q)+\varepsilon H((-q,q))\)
is strictly increasing, since \(A\) is strictly increasing and the second
term is nondecreasing. So at most one \(q\) solves \(F((-q,q))=A(k)\).

(i) Since \(A\) increases continuously from zero to one, \(q_1\) exists
uniquely. If \(q>q_1\), then
\(F((-q,q))\geq(1-\varepsilon)A(q)>(1-\varepsilon)A(q_1)=A(k)\), a
contradiction; so \(q\leq q_1\). Next
\(E_F\min(X^2,q^2)=(1-\varepsilon)B(q)+\varepsilon E_H\min(X^2,q^2)
\leq(1-\varepsilon)B(q)+\varepsilon q^2\). Both terms are nondecreasing
in \(q\), because \(\min(z^2,q^2)\) is, so the numerator is at most
\((1-\varepsilon)B(q_1)+\varepsilon q_1^2\). The denominator is
\(A(k)^2=(1-\varepsilon)^2A(q_1)^2\). Hence
\(V_1(F)\leq W(q_1,\varepsilon)\). If \(H\{|x|>q_1\}=1\), then
\(F((-q_1,q_1))=(1-\varepsilon)A(q_1)=A(k)\), so \(q=q_1\), and
\(E_H\min(X^2,q_1^2)=q_1^2\); equality follows.

(ii) Suppose \(A(k)\geq1-\varepsilon\). Since \(A(k)<1\),
\(A(k)-\varepsilon<1-\varepsilon\), so there is \(q_2\) with
\((1-\varepsilon)A(q)\geq A(k)-\varepsilon\) for \(q\geq q_2\).
For such \(q\) put \(h=\{A(k)-(1-\varepsilon)A(q)\}/\varepsilon\).
Then \(h\geq\{A(k)-(1-\varepsilon)\}/\varepsilon\geq0\) and \(h\leq1\).
Take \(M>q\) and
\(H=h\delta_0+\tfrac{1-h}2(\delta_{-M}+\delta_M)\), which is symmetric.
Then \(F((-q,q))=(1-\varepsilon)A(q)+\varepsilon h=A(k)\), so \(q\) is
the population cutoff, and
\(E_F\min(X^2,q^2)\geq\varepsilon(1-h)q^2\).
As \(q\to\infty\), \(\varepsilon(1-h)\to\varepsilon-A(k)+1-\varepsilon
=1-A(k)>0\). Therefore \(V_1(F)\geq\varepsilon(1-h)q^2/A(k)^2\to\infty\).

(iii) The equation \(A(k)=(1-\varepsilon)A(k_\varepsilon)\) has a unique
solution because \(A\) is a continuous increasing bijection onto
\((0,1)\). Then \(A(k)<1-\varepsilon\) and \(q_1=k_\varepsilon\), so part
(i) gives the worst case \(W(k_\varepsilon,\varepsilon)=1/I_*\) by
\cref{thm:p15-huber-least-favorable-normal}; this is the minimax value of
\cref{thm:p15-huber-location-minimax}.
\end{proof}
\end{annexitem}

\begin{annexitem}{annex:p15-w02b-p15-071-empirical-tools}{thm:p15-w02b-p15-071-empirical-tools}{Empirical tools for the count equation}
\begin{proof}
(i) For a nondecreasing \(g\) with values in
\([-1,1]\), \(g(x)=-1+\int_{-1}^1{\bf1}_{\{g(x)>s\}}\,ds\). Each set
\(G_s=\{x:g(x)>s\}\) is empty, the whole line, or a ray
\((c,\infty)\) or \([c,\infty)\), because \(g\) is nondecreasing.
Fubini's theorem gives
\(\int g\,dF_n-\int g\,dF=\int_{-1}^1\{F_n(G_s)-F(G_s)\}\,ds\).
For a ray \((c,\infty)\) the integrand is \(F(c)-F_n(c)\); for
\([c,\infty)\) it is \(F(c-)-F_n(c-)\), the limit of
\(F(y)-F_n(y)\) as \(y\uparrow c\). Both have absolute value at most
\(D_n\), so the integral is at most \(2D_n\). An interval's
probability is a difference of two values of \(F\) or of its left
limits, each pair within \(D_n\), which gives the second bound.

(ii) Let \(c(q)=\sup_tF([t-q,t+q])\), which is nonincreasing as
\(q\downarrow0\) and at least \(\max_xF\{x\}\); let \(L\) be its limit.
Suppose \(L>\max_xF\{x\}\). Choose \(R\) with
\(F(\{|x|>R\})<L/2\) and \(t_j\) with
\(F([t_j-1/j,t_j+1/j])>L-1/j\). For large \(j\) the interval meets
\([-R,R]\), so \(|t_j|\leq R+1\), and a subsequence converges to some
\(t^*\). For each \(r>0\), eventually
\([t_j-1/j,t_j+1/j]\subseteq[t^*-r,t^*+r]\), so
\(F([t^*-r,t^*+r])\geq L\). Letting \(r\downarrow0\),
\(F\{t^*\}\geq L\), a contradiction. Hence \(L=\max_xF\{x\}\); the
maximum exists because at most finitely many atoms exceed any positive
level.

(iii) Write \(r_i=|X_i|\) and \(\sigma_i\) for the sign of \(X_i\)
(zero when \(X_i=0\)). Since
\(\psi_q(x)=\operatorname{sign}(x)\min(|x|,q)\),
\[
 \psi_q(X_i)-\psi_\ell(X_i)=\sigma_i\{\min(r_i,q)-\min(r_i,\ell)\}
 =\pm\sigma_i\int_{\min(q,\ell)}^{\max(q,\ell)}{\bf1}_{\{r_i>u\}}\,du,
\]
with the plus sign when \(q\geq\ell\). Hence
\(|R_n(q)|\leq|q-\ell|\sup_{u>0}|W_n(u)|\), where
\(W_n(u)=n^{-1/2}\sum_i\sigma_i{\bf1}_{\{r_i>u\}}\).
For every fixed choice of signs \(\zeta\in\{-1,1\}^n\), the vector
\((\zeta_iX_i)\) has the law of \((X_i)\), by independence and
symmetry. Hence, conditionally on \((r_i)\), the signs \(\sigma_i\) of
the indices with \(r_i>0\) are independent and equal to \(\pm1\) with
probability one half each. Order the indices by decreasing \(r_i\),
breaking ties by index; every set \(\{i:r_i>u\}\) is an initial segment
of this order, so \(\sup_u|W_n(u)|\) is at most \(n^{-1/2}\) times the
largest absolute partial sum of the ordered signs. Conditionally on
\((r_i)\), \cref{thm:p15-kolmogorov-maximal} bounds the probability
that this exceeds \(x/\kappa\) by \(n\kappa^2/(nx^2)\). Integrating over
\((r_i)\) gives the bound.

(iv) Put \(f_0=F'(x_0)\), \(I_n=[x_0-M/\sqrt n,x_0+M/\sqrt n]\), fix an
integer \(r\geq1\) and the grid
\(y_j=x_0-M/\sqrt n+2jM/(r\sqrt n)\), \(0\leq j\leq r\). Let
\(\mathcal G_n\) be the family of the intervals \([y_i,y_j]\) and
\((y_i,y_j)\), \(i\leq j\), at most \(2(r+1)^2\) in number. Each
\(G\in\mathcal G_n\) has
\(\operatorname{Var}\{\sqrt n(F_n(G)-F(G))\}\leq F(G)\leq F(I_n)\), and
\(F(I_n)\to F\{x_0\}=0\) since differentiability gives continuity at
\(x_0\). By Chebyshev's inequality and a union bound,
\(\max_{G\in\mathcal G_n}\sqrt n|F_n(G)-F(G)|\to0\) in probability.
Differentiability gives \(F(y)=F(x_0)+f_0(y-x_0)+o(n^{-1/2})\)
uniformly for \(y\in I_n\), and the same for the left limits
\(F(y-)\). Hence every cell \([y_j,y_{j+1}]\) has
\(\sqrt nF([y_j,y_{j+1}])\leq2Mf_0/r+o(1)\). An interval
\(J\subseteq I_n\) lies between \(G_{\rm in}=(y_{i+1},y_{j-1})\) and
\(G_{\rm out}=[y_i,y_j]\) for suitable \(i,j\) (with
\(G_{\rm in}\) possibly empty), and \(G_{\rm out}\setminus G_{\rm in}\)
lies in two cells. So
\[
 F_n(J)-F(J)\leq F_n(G_{\rm out})-F(G_{\rm out})
             +F(G_{\rm out}\setminus G_{\rm in}),
\]
and similarly
\(F_n(J)-F(J)\geq F_n(G_{\rm in})-F(G_{\rm in})
-F(G_{\rm out}\setminus G_{\rm in})\). Therefore
\(\limsup_nP(\sup_J\sqrt n|F_n(J)-F(J)|>x)=0\) whenever
\(4Mf_0/r<x\); since \(r\) is arbitrary, this holds for every \(x>0\).

(v) Fix \(z\) and \(0<\eta<u/2\), and let \(G_n\) be the event
\(\{|U_n-u|\leq\eta,|V_n|\leq\eta\}\), whose probability tends to one.
On \(G_n\), \((Y_n+V_n)/U_n\leq z\) implies
\(Y_n\leq zU_n-V_n\leq zu+(|z|+1)\eta\), and is implied by
\(Y_n\leq zu-(|z|+1)\eta\). The normal limit has a continuous
distribution function, so
\[
 \Phi\!\left(\frac{zu-(|z|+1)\eta}\sigma\right)\leq
 \liminf_nP\!\left(\frac{Y_n+V_n}{U_n}\leq z\right)\leq
 \limsup_nP\!\left(\frac{Y_n+V_n}{U_n}\leq z\right)\leq
 \Phi\!\left(\frac{zu+(|z|+1)\eta}\sigma\right).
\]
Letting \(\eta\downarrow0\) gives the limit \(\Phi(zu/\sigma)\), the
distribution function of \(N(0,\sigma^2/u^2)\) at \(z\).
\end{proof}
\end{annexitem}

\begin{annexitem}{annex:p15-w02b-p15-071-prop1-consistency}{thm:p15-w02b-p15-071-prop1-consistency}{Existence and consistency of every Proposal 1 solution}
\begin{proof}
Write \(a=A(k)\in(0,1)\) and \(M_F=\max_xF\{x\}<a\).

\emph{Population facts.} For \(q<q'\),
\(F([-q',q'])-F([-q,q])\geq(1-\varepsilon)\{A(q')-A(q)\}>0\), so
\(q\mapsto F([-q,q])\) is strictly increasing; it tends to one as
\(q\to\infty\), so \(q_0<\infty\), and by part (ii) of
\cref{thm:p15-w02b-p15-071-empirical-tools} its limit as
\(q\downarrow0\) is at most \(M_F<a\), so \(q_0>0\). Right continuity
(the sets \([-q,q]\) decrease to \([-q_0,q_0]\)) gives
\(F([-q_0,q_0])=\alpha_0+2p_0\geq a\). For \(q<q_0\),
\(F([-q,q])<a\), and these sets increase to \((-q_0,q_0)\), so
\(\alpha_0\leq a\). Also \(\alpha_0\geq(1-\varepsilon)A(q_0)>0\).

\emph{Existence.} By the second bound of part (i) applied to single
points, \(\max_xF_n\{x\}\leq M_F+2D_n\), while \(m_n/n\to a\) and
\(D_n\to0\) almost surely by \cref{thm:p15-glivenko-cantelli}. So
almost surely \(E_n\) occurs for all large \(n\). On \(E_n\),
\(Q_n(t)>0\) for every \(t\). The \(m_n\)th smallest of the continuous
functions \(t\mapsto|x_i-t|\) is continuous, and \((q,u)\mapsto\psi_q(u)\)
is continuous, so \(g_n(t)=\sum_i\psi_{Q_n(t)}(x_i-t)\) is continuous.
For \(t<\min_ix_i\) every term is positive, and for \(t>\max_ix_i\)
every term is negative; the intermediate value theorem gives a zero,
and every zero \(T\) gives the solution \((T,Q_n(T)/k)\). The zero set
is closed and bounded below, so \(T_n\) exists. It is measurable:
\(T_n\leq c\) exactly when the minimum of \(g_n\) over
\([\min_ix_i-1,c]\) is at most zero, and that minimum equals the
infimum over the rationals of the interval together with \(c\).
Translating the sample translates every zero, and \(E_n\) and the
median are equivariant, so \(T^{(1)}\) is translation equivariant.

\emph{Location.} Let \(\lambda(t,q)=E_F\psi_1((X-t)/q)\). Since \(X\)
and \(-X\) have the same law and \(\psi_1\) is odd,
\[
 \lambda(t,q)=-\tfrac12E\{\psi_1(X/q+t/q)-\psi_1(X/q-t/q)\}.
\]
For \(t>0\) and real \(w\), \(\psi_1(w+\tau)-\psi_1(w-\tau)\) with
\(\tau=t/q\) is the length of \((w-\tau,w+\tau]\cap[-1,1]\), because
\(\psi_1(v)+1\) is the length of \((-\infty,v]\cap[-1,1]\). If
\(|w|\leq1\), the half of this interval on the side of zero lies in
\([-1,1]\) up to length \(\min(\tau,1)\). Hence
\[
 \lambda(t,q)\leq-\tfrac12\min(t/q,1)\,F([-q,q])\qquad(t>0),
\]
and \(\lambda(-t,q)=-\lambda(t,q)\).

By part (ii) of the tools there are \(q_{\min}>0\) and
\(0<\eta_0<\min(a-M_F,1-a)/2\) with \(F([t-q,t+q])<a-\eta_0\) for all
\(t\) whenever \(q<q_{\min}\). Choose \(R\) with
\(F((-R,R))>a+\eta_0\). Fix \(\delta>0\) and put
\(\gamma=\tfrac12\min\{\delta/(\delta+R),1\}(1-\varepsilon)A(q_{\min})>0\).
Almost surely, for all large \(n\), \(E_n\) holds,
\(2D_n+|m_n/n-a|<\eta_0\) and \(2D_n<\gamma\).
Let \((T,S)\) be a solution and \(q=kS=Q_n(T)\). By the definition of
the order statistic and part (i),
\[
 F([T-q,T+q])\geq\frac{m_n}n-2D_n>a-\eta_0,\qquad
 F((T-q,T+q))<\frac{m_n}n+2D_n<a+\eta_0 .
\]
The first gives \(q\geq q_{\min}\). If \(q>|T|+R\), the open window
would contain \((-R,R)\), contradicting the second; so
\(q\leq|T|+R\). Since \(\psi_q(u)=q\psi_1(u/q)\), the location equation
says \(n^{-1}\sum_i\psi_1((x_i-T)/q)=0\), and by part (i), applied to
the nondecreasing function \(x\mapsto\psi_1((x-T)/q)\),
\(|\lambda(T,q)|\leq2D_n<\gamma\). If \(T\geq\delta\), then
\(T/q\geq\delta/(\delta+R)\) and
\(F([-q,q])\geq(1-\varepsilon)A(q_{\min})\), so
\(\lambda(T,q)\leq-\gamma\), a contradiction; \(T\leq-\delta\) is
excluded by oddness. Hence every solution has \(|T|<\delta\).

\emph{Scale.} Fix \(0<\kappa<q_0\). Put
\[
 \theta=\min\bigl\{F([-q_0-\tfrac\kappa2,q_0+\tfrac\kappa2])-a,\;
 a-F([-q_0+\tfrac\kappa2,q_0-\tfrac\kappa2])\bigr\},
\]
which is positive: the first term is at least
\((1-\varepsilon)\{A(q_0+\kappa/2)-A(q_0)\}\), and the second is
positive because \(q_0-\kappa/2<q_0\). Take \(\delta<\kappa/2\) above
and \(n\) so large that also \(2D_n+|m_n/n-a|<\theta\). For a solution
with \(q\geq q_0+\kappa\), the open window contains
\([-q_0-\kappa/2,q_0+\kappa/2]\), so
\(a+\theta>F((T-q,T+q))\geq a+\theta\), which is impossible. For
\(q\leq q_0-\kappa\), the closed window lies in
\([-q_0+\kappa/2,q_0-\kappa/2]\), so
\(a-\theta<F([T-q,T+q])\leq a-\theta\), again impossible. So every
solution has \(|kS-q_0|<\kappa\), and \(\kappa\) was arbitrary.
\end{proof}
\end{annexitem}

\begin{annexitem}{annex:p15-w02b-p15-071-prop1-limit}{thm:p15-w02b-p15-071-prop1-limit}{Limiting risk of Proposal 1, including window atoms}
\begin{proof}
Write \(a=A(k)\), \(\Psi_n(t,q)=n^{-1}\sum_i\psi_q(X_i-t)\) and
\(N_n(s,q)=n^{-1}\#\{i:|X_i-s|<q\}\).

\emph{An exact identity.} For fixed \(x\) and \(q>0\), the function
\(s\mapsto\psi_q(x-s)\) is Lipschitz with derivative
\(-{\bf1}_{\{|x-s|<q\}}\) except at \(s=x\pm q\); the fundamental
theorem of calculus for Lipschitz functions gives
\(\psi_q(x-t)-\psi_q(x)=-\int_0^t{\bf1}_{\{|x-s|<q\}}\,ds\).
Averaging over the sample, every solution \((T,S)\) with \(q=kS\)
satisfies
\[
 T\,\overline N=\Psi_n(0,q),\qquad
 \overline N=\frac1T\int_0^TN_n(s,q)\,ds\quad(T\neq0).
\]
When \(T=0\) both sides vanish; put \(\overline N=c\) there, where
\(c\) is the constant of the case under discussion (\(\alpha_0\) in
(i)). Multiply by \(\sqrt n\) and write
\[
 \sqrt n\,\Psi_n(0,q)=Z_n+R_n(q),\qquad
 Z_n=n^{-1/2}\sum_i\psi_{q_0}(X_i),
\]
with \(R_n\) as in part (iii) of
\cref{thm:p15-w02b-p15-071-empirical-tools} for \(\ell=q_0\).

\emph{The two random terms.} The summands of \(Z_n\) are i.i.d.,
centered by symmetry, bounded by \(q_0\), with variance
\(\sigma_0^2>0\), because \(P(X\neq0)>0\). The row
\(\psi_{q_0}(X_i)/\sqrt n\) has its Lindeberg sums equal to zero once
\(q_0/\sqrt n\) is below the threshold, so
\cref{thm:p15-clt-lindeberg-feller} gives
\(Z_n\Rightarrow N(0,\sigma_0^2)\). Let \(\rho_n\) be the supremum of
\(|R_n(kS)|\) over all solutions. For \(\kappa>0\),
\(P(\rho_n\geq x)\leq P(\sup|kS-q_0|>\kappa)+\kappa^2/x^2\), and the
first term tends to zero by
\cref{thm:p15-w02b-p15-071-prop1-consistency}; so \(\rho_n\to0\) in
probability.

\emph{Bracketing the average count.} If \(|T|+|q-q_0|\leq\eta\) and
\(s\) lies between \(0\) and \(T\), then
\((-q_0+\eta,q_0-\eta)\subseteq(s-q,s+q)\subseteq(-q_0-\eta,q_0+\eta)\).
Hence
\[
 F_n((-q_0+\eta,q_0-\eta))\leq\overline N\leq
 F_n((-q_0-\eta,q_0+\eta)),
\]
whenever \(T\neq0\). As \(\eta\downarrow0\) the population values of the
two bounds tend to \(\alpha_0\) and \(\alpha_0+2p_0\), and by part (i) of
the tools and \cref{thm:p15-glivenko-cantelli} the empirical values are
within \(2D_n\to0\) of them.

(i) Fix \(0<\delta'<\alpha_0/2\), and choose \(\eta\) with
\(F((-q_0+\eta,q_0-\eta))>\alpha_0-\delta'/2\). On the event
\(G_n=\{\rho_n\leq\delta'\}\cap\{2D_n<\delta'/2\}\cap
\{\text{every solution has }|T|+|kS-q_0|\leq\eta\}\), whose probability
tends to one, every solution has \(\overline N\geq\alpha_0-\delta'\)
and \(\sqrt n|T|\leq(|Z_n|+\delta')/(\alpha_0-\delta')\). For \(C>0\),
\[
 E\min(nT_n'^2,C^2)\leq
 E\min\!\left(\frac{(|Z_n|+\delta')^2}{(\alpha_0-\delta')^2},C^2\right)
 +C^2P(G_n^c).
\]
The function inside the first expectation is bounded and continuous in
\(Z_n\), so by \cref{thm:p15-weak-bounded-continuous} the right side
tends to \(E\min((|Z|+\delta')^2/(\alpha_0-\delta')^2,C^2)\), with
\(Z\sim N(0,\sigma_0^2)\). Let \(\delta'\downarrow0\) by bounded
convergence, then \(C\to\infty\) by monotone convergence; the bound is
\(EZ^2/\alpha_0^2=\sigma_0^2/\alpha_0^2\).

\emph{Exact limits.} In (ii)--(iv) it is shown below that
\(\overline N\to c\) in probability, uniformly over solutions, for a
constant \(c>0\). Then \(\sqrt nT_n'=(Z_n+R_n(kS_n'))/\overline N\) with
\(|R_n(kS_n')|\leq\rho_n\), and part (v) of the tools gives
\(\sqrt nT_n'\Rightarrow N(0,\sigma_0^2/c^2)\). The limiting truncated
risk is then the limit's second moment, by the preliminary fact at the
head of this section.

(ii) If \(p_0=0\), both bracketing bounds tend to \(\alpha_0\), so
\(c=\alpha_0\); and \(\alpha_0\leq a\leq\alpha_0+2p_0=\alpha_0\) gives
\(\alpha_0=a\). Then \(F((-q_0,q_0))=a\), so \(q_0\) is the cutoff of
\cref{thm:p15-w02b-p15-071-prop1-population}, and the variance is
\(\sigma_0^2/a^2=V_1(F)\).

(iii) Let \(\theta=\min(a-\alpha_0,\alpha_0+p_0-a)/4>0\), and choose
\(0<\delta<q_0/2\) with \(F((-q_0,-q_0+2\delta])<\theta\); by symmetry
also \(F([q_0-2\delta,q_0))<\theta\). Almost surely, for all large
\(n\), every solution has \(|T|<\delta\), and
\(4D_n+|m_n/n-a|<\theta\). Let \(0\leq T<\delta\) and
\(q'=q_0-T>0\). The closed window \([2T-q_0,q_0]\) contains
\((-q_0+2\delta,q_0)\) and the atom \(q_0\), so its empirical mass is at
least \(\alpha_0-\theta+p_0-4D_n>a+2\theta>m_n/n\). The open window
\((2T-q_0,q_0)\) lies in \((-q_0,q_0)\), so its empirical mass is at
most \(\alpha_0+2D_n<a-3\theta<m_n/n\). Hence \(Q_n(T)=q'\). The case
\(-\delta<T<0\) is the mirror image, with the atom \(-q_0\). So every
solution has \(kS=q_0-|T|\). For \(s\) between \(0\) and \(T\), the
window \((s-q_0+|T|,s+q_0-|T|)\) lies in \((-q_0,q_0)\) and contains
\((-q_0+2|T|,q_0-2|T|)\). So
\(F_n((-q_0+2|T|,q_0-2|T|))\leq\overline N\leq F_n((-q_0,q_0))\), and
both bounds tend to \(\alpha_0\) in probability, uniformly over
solutions, because \(|T|\to0\). Thus \(c=\alpha_0\).

(iv) Let \(\theta=\min(a-\alpha_0-p_0,\alpha_0+2p_0-a)/4>0\) and choose
\(\delta>0\) with \(F((q_0,q_0+2\delta))<\theta\). For large \(n\) as
before, let \(0\leq T<\delta\) and \(q'=q_0+T\). The closed window
\([-q_0,q_0+2T]\) contains \([-q_0,q_0]\), with empirical mass at least
\(\alpha_0+2p_0-2D_n>m_n/n\). The open window \((-q_0,q_0+2T)\) has
empirical mass at most \(\alpha_0+p_0+\theta+2D_n<m_n/n\). Hence
\(Q_n(T)=q_0+T\), and the case \(T<0\) is the mirror image. For \(s\)
strictly between \(0\) and \(T\), the window
\((s-q_0-|T|,s+q_0+|T|)\) contains \([-q_0,q_0]\) and lies in
\((-q_0-2|T|,q_0+2|T|)\), so
\(F_n([-q_0,q_0])\leq\overline N\leq F_n((-q_0-2|T|,q_0+2|T|))\), and
both bounds tend to \(\alpha_0+2p_0\). Thus \(c=\alpha_0+2p_0\).

(v) Let \(f=F'(q_0)\). By symmetry \(F\) is differentiable at
\(-q_0\) with the same derivative, and continuity gives \(p_0=0\).
Since \(F-(1-\varepsilon)\Phi\) is nondecreasing,
\(f\geq(1-\varepsilon)\phi(q_0)>0\); and \(F([-q_0,q_0])=a\) by (ii).
By (ii), \(\sqrt nT_n'\) converges in law, so for \(\eta>0\) there is
\(M\) with \(P(\sqrt n|T_n'|>M)<\eta\) for all large \(n\).

Fix \(z\), put \(q'=q_0+z/\sqrt n\) and \(q=Q_n(T_n')\). If \(q>q'\)
then \([T_n'-q',T_n'+q']\subseteq(T_n'-q,T_n'+q)\), whose count is
below \(m_n\); if \(q\leq q'\) then \([T_n'-q',T_n'+q']\) contains
\([T_n'-q,T_n'+q]\), whose count is at least \(m_n\). Hence
\[
 \{\sqrt n(q-q_0)>z\}=\{F_n([T_n'-q',T_n'+q'])<m_n/n\}.
\]
Write \(F_n([T_n'-q',T_n'+q'])-a=\Delta_1+\Delta_2+\Delta_3\), where
\(\Delta_1=F_n([T_n'-q',T_n'+q'])-F_n([-q',q'])\),
\(\Delta_2=(F_n-F)([-q',q'])\) and \(\Delta_3=F([-q',q'])-a\).
The set difference in \(\Delta_1\) consists of two intervals inside
\([\pm q_0-(M+|z|)/\sqrt n,\pm q_0+(M+|z|)/\sqrt n]\) when
\(\sqrt n|T_n'|\leq M\). Their population masses are
\(f|T_n'|+o(n^{-1/2})\) each, entering with opposite signs, and part
(iv) of the tools controls their empirical deviations; so
\(\sqrt n\Delta_1\to0\) in probability. Part (iv) also gives
\(\sqrt n\{\Delta_2-(F_n-F)([-q_0,q_0])\}\to0\), and
\(\sqrt n(F_n-F)([-q_0,q_0])\Rightarrow N(0,a(1-a))\) by
\cref{thm:p15-clt-lindeberg-feller} for the bounded indicators.
Differentiability gives \(\sqrt n\Delta_3\to2fz\), and
\(\sqrt n|m_n/n-a|\leq n^{-1/2}\). The event above is therefore
\(\{Y_n+V_n<-2fz\}\) with \(Y_n\Rightarrow N(0,a(1-a))\) and
\(V_n\to0\) in probability, and part (v) of the tools (with
\(U_n=1\)) gives
\(P(\sqrt n(q-q_0)>z)\to\Phi(-2fz/\sqrt{a(1-a)})\). So
\(\sqrt n(kS_n'-q_0)\Rightarrow N(0,a(1-a)/(4f^2))\).
By the mean-value theorem
\(\sqrt n\{\log S_n'-\log(q_0/k)\}=\sqrt n(kS_n'-q_0)/\xi_n\) with
\(\xi_n\) between \(kS_n'\) and \(q_0\), so \(\xi_n\to q_0\) almost
surely, and part (v) of the tools gives the stated limit.
\end{proof}
\end{annexitem}

\begin{annexitem}{annex:p15-w02b-p15-071-prop1-minimax}{thm:p15-w02b-p15-071-prop1-minimax}{Proposal 1's asymptotic minimax standing}
\begin{proof}
Write \(a=A(k)\).

(i) Since \(F([-q_1,q_1])\geq(1-\varepsilon)A(q_1)=a\), \(q_0\leq q_1\).
Further \(\sigma_0^2=(1-\varepsilon)B(q_0)+\varepsilon E_H\min(X^2,q_0^2)
\leq(1-\varepsilon)B(q_0)+\varepsilon q_0^2\), using
\(B(q)=E\min(Z^2,q^2)\) from the proof of
\cref{thm:p15-w02b-p15-071-bias-variance}(iii), and
\(\alpha_0\geq(1-\varepsilon)A(q_0)\). Part (i) of
\cref{thm:p15-w02b-p15-071-prop1-limit} gives the bound
\(\sigma_0^2/\alpha_0^2\leq V^*(q_0,\varepsilon)\).

(ii) For regular \(F\), part (ii) of
\cref{thm:p15-w02b-p15-071-prop1-limit} gives the risk \(V_1(F)\) with
\(q_0\) the population cutoff, and
\cref{thm:p15-w02b-p15-071-prop1-population}(i) gives
\(V_1(F)\leq V^*(q_1,\varepsilon)\). If \(H\) is supported on
\(\{|x|>q_1\}\), then \(F([-q,q])=(1-\varepsilon)A(q)\) for
\(q\leq q_1\), so \(q_0=q_1\) and \(p_0=\varepsilon H\{q_1\}=0\); the
same theorem gives equality.

(iii) The equation for \(k\) gives \((1-\varepsilon)A(q_1)=a=
(1-\varepsilon)A(k_\varepsilon)\), so \(q_1=k_\varepsilon\), and
\(V^*(k_\varepsilon,\varepsilon)=1/I_*\) by
\cref{thm:p15-huber-least-favorable-normal}.
Since \((1-\varepsilon)A\) increases continuously from zero, \(q_*\)
is well defined and \(q_*<k_\varepsilon\). Fix
\(q\in(q_*,k_\varepsilon)\). Then \((1-\varepsilon)A(q)<a\) and
\(a-(1-\varepsilon)A(q)<\min(a,\varepsilon/2)\): the second inequality
is \((1-\varepsilon)A(q)>a-\varepsilon/2\) when \(a>\varepsilon/2\),
and holds because \((1-\varepsilon)A(q)>0\) otherwise. Choose \(h\) with
\(a-(1-\varepsilon)A(q)<\varepsilon h<\min(a,\varepsilon/2)\), and let
\[
 H=h(\delta_q+\delta_{-q})+(1-2h)U,
\]
with \(U\) uniform on \([-q-2,-q-1]\cup[q+1,q+2]\). It is symmetric;
the only atoms of \(F\) are at \(\pm q\), of mass \(\varepsilon h<a\),
so (I) holds. For \(q'<q\), \(F([-q',q'])=(1-\varepsilon)A(q')<a\), and
\(F([-q,q])=(1-\varepsilon)A(q)+2\varepsilon h>a\); so \(q_0=q\),
\(\alpha_0=(1-\varepsilon)A(q)<a<\alpha_0+\varepsilon h=\alpha_0+p_0\),
which is configuration (iii). All of \(H\) lies in \(\{|x|\geq q\}\),
so \(\sigma_0^2=(1-\varepsilon)B(q)+\varepsilon q^2\), and the limiting
truncated risk is \(\sigma_0^2/\alpha_0^2=V^*(q,\varepsilon)\).

For the strict inequality, the least favorable law \(F_*\) of
\cref{thm:p15-huber-least-favorable-normal} lies in
\(\mathcal C_\varepsilon\) and has no atoms, so by
\cref{thm:p15-huber-normal-variance}
\(V^*(q,\varepsilon)\geq E_{F_*}\psi_q^2/F_*((-q,q))^2\). Here
\(E_{F_*}\psi_q=0\), \(0<E_{F_*}\psi_q^2<\infty\), and
\(F_*((-q,q))=-\frac d{dt}E_{F_*}\psi_q(X-t)|_{t=0}\) by the
derivative formula in the proof of
\cref{thm:p15-w02b-p15-071-bias-variance}(iii). The same theorem
gives \(E_{F_*}\psi_q^2/F_*((-q,q))^2\geq1/I_*\), with equality only if
\(\psi_q\) is a multiple of \(\psi_{k_\varepsilon}\) almost everywhere
for the positive density \(f_*\). On \((0,q)\) the two scores agree, so
the multiple would be one; but they differ on
\((q,k_\varepsilon)\), a set of positive \(F_*\)-probability.
Hence \(V^*(q,\varepsilon)>1/I_*\).
The sequence \(T^{(1)}\) is translation equivariant, the laws just
built belong to \(\mathcal C_\varepsilon\), and their limiting
truncated risks exist, so
\(\mathcal R_\varepsilon(T^{(1)})\geq V^*(q,\varepsilon)\) for each
such \(q\).

(iv) Every symmetric \(H\) satisfies (I). By (i) and
\(\liminf\leq\limsup\),
\(\mathcal R_\varepsilon(T^{(1)})\leq\sup_{0<q\leq q_1}V^*(q,\varepsilon)\).
The function \(V^*(\cdot,\varepsilon)\) is continuous on \((0,q_1]\).
As \(q\downarrow0\), \(B(q)/q^2=E\min(Z^2/q^2,1)\to1\) by dominated
convergence, and \(A(q)/q\to2\phi(0)\), so
\(V^*(q,\varepsilon)\to\{(1-\varepsilon)+\varepsilon\}/\{(1-\varepsilon)^2
4\phi(0)^2\}=\pi/\{2(1-\varepsilon)^2\}\). The supremum is therefore
finite.

The numerical table is an evaluation of these equations, and the
value \(\sigma_0^2/A(k)^2\) quoted for comparison uses
\(\sigma_0^2=V^*(q,\varepsilon)\alpha_0^2\). For the wrong-\(k\)
remark, the laws in the proof of
\cref{thm:p15-w02b-p15-071-prop1-population}(ii) have atoms only at
\(0\) and \(\pm M\) with \(M>q\), so they are regular, and their atom
masses are at most \(\varepsilon<a\).
\end{proof}
\end{annexitem}

\begin{annexitem}{annex:p15-w02b-p15-071-scale-compare}{thm:p15-w02b-p15-071-scale-compare}{The scale accuracy comparison at the normal law}
\begin{proof}
(i) For \(F=\Phi\), (I) holds since \(\Phi\) has no atoms,
\(q_0=k\) because \(\Phi([-k,k])=A(k)\), and \(\Phi\) is differentiable
at \(k\) with derivative \(\phi(k)\). Part (v) of
\cref{thm:p15-w02b-p15-071-prop1-limit}, with \(\log(q_0/k)=0\), gives
the limit.

(ii) Under \(\Phi\), \(\Phi\{0\}=0<1-B(k)/k^2\), so
\cref{thm:p15-huber-joint-identified} applies, and \(S_0=1\) because
\(E\psi(k,Z)^2=E\min(Z^2,k^2)=B(k)\). With \(U=Z\), no atoms occur, so
\cref{thm:p15-huber-joint-law-and-variance} gives
\(\sqrt n(S_n-1)\Rightarrow N(0,\eta^2/C^2)\), where
\(C=E[Z^2{\bf1}_{\{|Z|<k\}}]\) and
\(\eta^2=\operatorname{Var}(\psi(k,Z)^2)/4\). Integration by parts gives
\(\int_{-k}^kz^2\phi=A(k)-2k\phi(k)=C(k)\) and
\(\int_{-k}^kz^4\phi=-2k^3\phi(k)+3C(k)\), so
\(E\psi(k,Z)^4=3A(k)-6k\phi(k)-2k^3\phi(k)+k^4(1-A(k))=M_4(k)\) and
\(\eta^2/C^2=v_2(k)\). The mean-value theorem, with the intermediate
point tending to one, and part (v) of
\cref{thm:p15-w02b-p15-071-empirical-tools} transfer the limit to
\(\log S_n\).

(iii) At \(k=1\), \(1-\Phi(1)=(1-A(1))/2\) gives \(B(1)=1-2\phi(1)\) and
\(M_4(1)=2A(1)+1-8\phi(1)\), so with \(u=A(1)-2\phi(1)\),
\[
 v_2(1)=\frac{u-2\phi(1)^2}{2u^2},\qquad
 v_1(1)=\frac{A(1)(1-A(1))}{4\phi(1)^2}.
\]
At \(k=2\), \(B(2)=4-3A(2)-4\phi(2)\) and \(M_4(2)=16-13A(2)-28\phi(2)\).
The computed values are \(\phi(1)\simeq0.241971\),
\(A(1)\simeq0.682689\), \(\phi(2)\simeq0.053991\),
\(A(2)\simeq0.954500\); use only the enclosures
\(\phi(1)\in[0.2419,0.2420]\), \(A(1)\in[0.6826,0.6827]\),
\(\phi(2)\in[0.0539,0.0541]\), \(A(2)\in[0.9544,0.9546]\).
The function \(A(1-A)/(4k^2\phi^2)\) decreases in \(A\) for \(A>1/2\)
and in \(\phi\). Hence \(v_1(1)\leq0.6826\cdot0.3174/(4\cdot0.2419^2)
<0.9257\) and
\(v_1(2)\geq0.9546\cdot0.0454/(16\cdot0.0541^2)>0.9254\).
For \(v_2(1)\), \(u\in[0.1986,0.1989]\), so
\(v_2(1)\geq(0.1986-2\cdot0.2420^2)/(2\cdot0.1989^2)>1.029\).
For \(v_2(2)\), put \(w=B(2)\in[0.9198,0.9212]\) and
\(u=A(2)-4\phi(2)\geq0.7380\); the numerator
\(M_4(2)-w^2\leq16-13\cdot0.9544-28\cdot0.0539-0.9198^2<1.2376\), so
\(v_2(2)<1.2376/(4\cdot0.7380^2)<0.5681\). These give (iii).
The tables and the crossing are evaluations of the displayed formulas.
\end{proof}
\end{annexitem}

\begin{annexitem}{annex:p15-w02b-p15-071-pair-median-risk}{thm:p15-w02b-p15-071-pair-median-risk}{The source bound for arbitrary symmetric contamination}
\begin{proof}
Write \(N=\binom n2\), \(Z_{ij}=(X_i+X_j)/2\), and
\(K(t)=P(X_1+X_2\leq2t)\).

\emph{A variance decomposition.} Let \(h\) be a symmetric measurable
kernel with values in \([0,1]\), \(\theta=Eh(X_1,X_2)\),
\(g(x)=Eh(x,X_2)\), \(\zeta=\operatorname{Var}g(X_1)\), and
\(U=N^{-1}\sum_{i<j}h(X_i,X_j)\). Put
\(\tilde h(x,y)=h(x,y)-g(x)-g(y)+\theta\). Counting each index in
\(n-1\) pairs,
\[
 U-\theta=\frac2n\sum_i\{g(X_i)-\theta\}+R,\qquad
 R=N^{-1}\sum_{i<j}\tilde h(X_i,X_j).
\]
Since \(E[\tilde h(x,X_2)]=0\) for every \(x\), two terms of \(R\)
with distinct index pairs are uncorrelated: if the pairs are disjoint
this follows from independence, and if they share one index, condition
on the shared and one unshared variable and integrate the other.
Also \(E\tilde h^2=E(h-\theta)^2-2\zeta\leq E(h-\theta)^2\leq1/4\).
Hence \(ER^2\leq1/(4N)\), and the projection term has variance
\(4\zeta/n\leq1/n\).

\emph{Symmetry.} Since \((X_i)\) and \((-X_i)\) have the same law and
negating all observations negates every \(Z_{ij}\), hence every order
statistic and the midpoint median, \(T_n\) and \(-T_n\) have the same
law. In particular \(P(T_n>t)=P(T_n<-t)\).

\emph{Order statistics.} The midpoint median is the average of
\(Z_{(\lceil N/2\rceil)}\leq Z_{(\lfloor N/2\rfloor+1)}\).
If \(T_n>t\), then \(Z_{(\lfloor N/2\rfloor+1)}>t\), so at most
\(\lfloor N/2\rfloor\) of the \(Z_{ij}\) are at most \(t\). Hence
\(P(T_n>t)\leq P(\widehat U_n(t)\leq1/2)\), where
\(\widehat U_n(t)=N^{-1}\#\{i<j:Z_{ij}\leq t\}\) is the statistic with
kernel \(h_t(x,y)={\bf1}_{\{x+y\leq2t\}}\), mean \(K(t)\) and projection
\(g_t(x)=F(2t-x)\).

(i) Let \(\varepsilon>0\) and \(H\{a\}=p>0\). By symmetry \(H\{-a\}=p\),
so \(\pi_0=P(X_1+X_2=0)\geq P(X_1=a,X_2=-a)\geq(\varepsilon p)^2>0\).
By symmetry \(P(X_1+X_2<0)=P(X_1+X_2>0)=(1-\pi_0)/2\).
Apply the decomposition to \(h_0\) and to
\({\bf1}_{\{x+y<0\}}\): their statistics have means \((1+\pi_0)/2\) and
\((1-\pi_0)/2\) and variances at most \(1/n+1/(4N)\). By Chebyshev's
inequality, with probability tending to one,
\(\#\{Z_{ij}<0\}<N/2-1\) and \(\#\{Z_{ij}\leq0\}>N/2+1\).
The first gives \(Z_{(\lceil N/2\rceil)}\geq0\), and the second gives
\(Z_{(\lfloor N/2\rfloor+1)}\leq0\). Thus both are zero and \(T_n=0\).
Then \(E\min(nT_n^2,C^2)\leq C^2P(T_n\neq0)\to0\) for every \(C\).

(ii) Let \(\varepsilon=0\) or let \(H\) have no atom. Then \(F\) is continuous and
\(P(X_1+X_2=0)=0\), so \(K(0)=1/2\). With \(Z_1,Z_2\sim\Phi\) and
\(Y_1,Y_2\sim H\), all independent, conditioning on the component labels
gives
\[
 K(t)=(1-\varepsilon)^2\Phi(\sqrt2\,t)
      +2\varepsilon(1-\varepsilon)L(t)
      +\varepsilon^2P(Y_1+Y_2\leq2t),
\]
where \(L(t)=\int\Phi(2t-y)\,H(dy)\),
because \(Z_1+Z_2\) is normal with variance two. Each term is
nondecreasing. By dominated convergence \(L'(0)=2\int\phi(y)\,H(dy)\),
using that \(\phi\) is even and bounded. Dropping the third increment,
for \(t>0\),
\[
 K(t)-\tfrac12\geq(1-\varepsilon)^2\{\Phi(\sqrt2\,t)-\tfrac12\}
 +2\varepsilon(1-\varepsilon)\{L(t)-L(0)\}.
\]
Fix \(z>0\) and \(t_n=z/\sqrt n\). Then
\(\liminf_n\sqrt n\{K(t_n)-1/2\}\geq
z\{(1-\varepsilon)^2\sqrt2\,\phi(0)+4\varepsilon(1-\varepsilon)\int\phi\,dH\}
=\kappa z\), since \(\sqrt2\,\phi(0)=1/\sqrt\pi\).

By right continuity \(g_{t_n}(x)=F(2t_n-x)\to F(-x)\) for every \(x\),
and bounded convergence gives
\(\zeta_n=\operatorname{Var}g_{t_n}(X)\to\operatorname{Var}F(-X)\).
Since \(F\) is continuous and symmetric, \(F(-X)=1-F(X)\), and
\(F(X)\) is uniform on \((0,1)\): for \(0<u<1\) let \(x_u\) be the
largest \(x\) with \(F(x)=u\), which exists by continuity; then
\(F(X)\leq u\) exactly when \(X\leq x_u\), with probability
\(F(x_u)=u\). Hence \(\zeta_n\to1/12\).
The row \(Y_{n,i}=2n^{-1/2}\{g_{t_n}(X_i)-K(t_n)\}\), \(i\leq n\),
consists of independent centered entries bounded by \(2n^{-1/2}\), with
total variance \(4\zeta_n\to1/3\); the Lindeberg sums vanish once
\(2n^{-1/2}\) is below the threshold. By
\cref{thm:p15-clt-lindeberg-feller}, \(\sum_iY_{n,i}\Rightarrow N(0,1/3)\),
and \(E(\sqrt nR)^2\leq n/(4N)\to0\). Therefore
\(W_n=\sqrt n\{\widehat U_n(t_n)-K(t_n)\}\Rightarrow N(0,1/3)\).

For \(0<\delta<\kappa\), eventually
\(\sqrt n\{1/2-K(t_n)\}\leq-(\kappa-\delta)z\), so
\[
 P(\sqrt nT_n>z)\leq P(W_n\leq-(\kappa-\delta)z)
 \longrightarrow\Phi(-\sqrt3(\kappa-\delta)z),
\]
the limit law having a continuous distribution function. Letting
\(\delta\downarrow0\) gives the tail bound, and symmetry gives the
lower tail.
For the risk, \(E\min(nT_n^2,C^2)=\int_0^{C^2}P(nT_n^2>v)\,dv\).
The integrand is at most one on a finite interval, so the reverse Fatou
lemma gives
\[
 \begin{split}
 \limsup_nE\min(nT_n^2,C^2)&\leq\int_0^{C^2}2\Phi(-\sqrt3\,\kappa\sqrt v)\,dv\\
 &\leq\int_0^\infty P\!\left(\frac{Z^2}{3\kappa^2}>v\right)dv
 =\frac1{3\kappa^2},
 \end{split}
\]
with \(Z\) standard normal.

\emph{The bound, strictness and sharpness.} Since
\(\kappa\geq(1-\varepsilon)^2/\sqrt\pi\),
\(1/(3\kappa^2)\leq\pi/\{3(1-\varepsilon)^4\}\); case (i) gives zero.
For \(0<\varepsilon<1\), \(\int\phi\,dH>0\) because \(\phi>0\), so
\(\kappa>(1-\varepsilon)^2/\sqrt\pi\) and the inequality is strict in
case (ii); case (i) is strict trivially.
For symmetric contaminants with square-integrable densities,
\cref{thm:p15-huber-pair-median} gives a normal limit, whose variance is
the risk by the preliminary fact, and shows that their supremum is the
bound; hence the supremum over all symmetric \(H\) is the bound.
If \(\sqrt nT_n\) converges in law, the preliminary fact identifies the
risk with the limit's second moment.
\end{proof}
\end{annexitem}
